\hypertarget{supplementary-information-first-order-conflation-and-recurrent-fragmentation-generate-a-cognitive-near-singularity}{%
\section{Supplementary Information: First-order conflation and recurrent
fragmentation generate a cognitive
near-singularity}\label{supplementary-information-first-order-conflation-and-recurrent-fragmentation-generate-a-cognitive-near-singularity}}

\textbf{Andy E. Williams}

\emph{Manuscript version 16 \textbar{} 17 July 2026}

\hypertarget{contents-and-status}{%
\subsection{Contents and status}\label{contents-and-status}}

This Supplement is the standalone formal, computational and status
record for Article v15. It restores every load-bearing definition and
premise used by the Article, separates semantic registration from proof
checking, retains the executed finite realization from the earlier
computational package, adds the executed direct recurrence experiment
for this version, and records all unexecuted gates explicitly. No
unpublished artifact is required to recover the formal content.

\begin{longtable}[]{@{}
  >{\raggedright\arraybackslash}p{(\columnwidth - 4\tabcolsep) * \real{0.3333}}
  >{\raggedright\arraybackslash}p{(\columnwidth - 4\tabcolsep) * \real{0.3333}}
  >{\raggedright\arraybackslash}p{(\columnwidth - 4\tabcolsep) * \real{0.3333}}@{}}
\toprule\noalign{}
\begin{minipage}[b]{\linewidth}\raggedright
Package
\end{minipage} & \begin{minipage}[b]{\linewidth}\raggedright
Content
\end{minipage} & \begin{minipage}[b]{\linewidth}\raggedright
v15 status
\end{minipage} \\
\midrule\noalign{}
\endhead
\bottomrule\noalign{}
\endlastfoot
A & Core FMI-CNS mathematics & Prose proofs and countermodels supplied;
kernel formalization pending \\
B & Endogenous semantic geometry & Definitions and theorem targets
supplied; physical mapping remains a hypothesis \\
C & Lean verification artifact & Module and declaration manifest
supplied; pinned source package passes a static no-sorry /
no-custom-axiom scan; kernel build and axiom report pending \\
D & String-specific extension & Candidate geometric representation;
comparative utility experiment pending \\
E & External and propagation evidence & Positive v8 Gate Zero
reproduction; canonical recurrence rerun; bounded same-model receiver
study and bounded branch/hazard and mutation experiments executed;
cross-model receiver and external trace studies open \\
\end{longtable}

The package status tuple is Status(v16) = (conditionally closed,
conditionally closed, partially closed, open, closed-bounded), where
Package A/B carry prose/source proofs, Package C (Lean) is
source-complete and static-audit-clean with the kernel build blocked by
the environment's egress policy, Package D (string action) remains open,
and Package E now carries a positive v8 Gate Zero reproduction, a
canonical recurrence rerun, a bounded same-model receiver study and
bounded Great Filter branch/hazard and mutation experiments.

For the five submission gates addressed in the companion open-gate
closure record, the gate status tuple is

\begin{quote}
Status\_gates = (\emph{s}\_Lean, \emph{s}\_G0, \emph{s}\_Rec,
\emph{s}\_Recv, \emph{s}\_Ext) = (partial, closed-positive, closed,
closed-bounded, partial).
\end{quote}

The Lean kernel build is partially closed (source-complete, static-audit
clean; toolchain download egress-blocked); the v16 Gate Zero audit is
closed-positive (the v8 lineage was supplied and reproduced
bit-for-bit); the direct recurrence experiment is closed (canonical
common-latent-world rerun); the same-model receiver study is
closed-bounded (executed with isolated agents); the external mappings
are partially closed (registration-closed; no physical instantiation
promoted). Formal coherence within stated premises and
submission-carrier readiness are separate verdicts. No aggregate label
may erase which gate is open.

\hypertarget{supplementary-note-1.-typed-metric-functional-state-spaces-and-endogenous-semantics}{%
\subsection{Supplementary Note 1. Typed metric functional state spaces
and endogenous
semantics}\label{supplementary-note-1.-typed-metric-functional-state-spaces-and-endogenous-semantics}}

\emph{Local map.} This Note defines the mathematical object on which
every later theorem operates. Its payoff is a representation of semantic
identity that is constituted by metric position, typed interactions and
admissible relational configuration rather than by a prose label.

\hypertarget{definition-s1-typed-metric-fss}{%
\subsubsection{Definition S1: typed metric
FSS}\label{definition-s1-typed-metric-fss}}

A typed metric functional state space is \emph{X} = (\emph{V}, \emph{E},
τ\_V, τ\_E, \emph{g}, π, β, \emph{A}), where \emph{V} is the set of
represented functional states, \emph{E} the set of typed interaction or
transition instances, τ\_V and τ\_E the node and interaction typing
maps, \emph{g} a metric or metric family, π an optional coordinate or
embedding representation, β the boundary record and \emph{A} the family
of admissible transformations and trajectories.

A boundary record may contain domain and codomain restrictions,
certificates, provenance, residuals, successor obligations, fitness
relations, continuation conditions, version identities and target scope.
The metric may be primitive, interaction-induced or jointly defined.
Every theorem must state which case it uses.

\hypertarget{definition-s2-nodes-processes-and-functions}{%
\subsubsection{Definition S2: nodes, processes and
functions}\label{definition-s2-nodes-processes-and-functions}}

A represented concept or functional state is a node \emph{v} ∈ \emph{V}.
A reasoning or physical interaction is a typed edge, hyperedge, path or
process subgraph. Definitions, claims, functions, proofs, experiments,
datasets and manuscript passages are semantic roles borne by nodes and
processes rather than additional primitive represented-object kinds.

A represented function has an intensional node carrying identity,
domain, codomain and contract, and an extensional family of process
instances carrying executions. An \emph{instantiates} relation connects
the contract node to those instances.

\hypertarget{definition-s3-semantic-profile-and-identity}{%
\subsubsection{Definition S3: semantic profile and
identity}\label{definition-s3-semantic-profile-and-identity}}

For \emph{v} ∈ \emph{V}, Prof\_X(\emph{v}) = (\emph{X}, \emph{v}) and
Sem\_X(\emph{v}) = {[}(\emph{X}, \emph{v}){]}\_≅\emph{A}. An admissible
pointed isomorphism \emph{F} : (\emph{X}, \emph{v}) → (\emph{Y},
\emph{w}) must satisfy \emph{F}(\emph{v}) = \emph{w}, preserve node and
interaction types, incidence, orientation and the declared metric
relations, preserve boundary, certificate, warrant and continuation
records, and intertwine the relevant transition and participant-fitness
relations. Approximate isomorphisms must expose the tolerance and
residual rather than silently replacing equality.

\hypertarget{theorem-s1-semantic-equivalence}{%
\subsubsection{Theorem S1: semantic
equivalence}\label{theorem-s1-semantic-equivalence}}

If \emph{F} : (\emph{X}, \emph{v}) → (\emph{Y}, \emph{w}) is an
admissible target-preserving pointed isomorphism, then every invariant
target proposition expressible in the registered language has the same
truth value at \emph{v} and \emph{w}.

\emph{Proof.} Let the registered language be generated from atomic
predicates on node type, edge type, metric relation, incidence,
boundary, warrant, transition, fitness and continuation structure by
Boolean connectives and bounded quantification over the represented FSS.
Admissibility preserves every atomic predicate. The result follows by
induction on formula construction. For a quantified formula, the
bijection induced by \emph{F} transports witnesses and counterexamples.
Therefore truth is invariant. QED.

\hypertarget{definition-s4-finite-target-complete-semantic-kernel}{%
\subsubsection{Definition S4: finite target-complete semantic
kernel}\label{definition-s4-finite-target-complete-semantic-kernel}}

For target family \emph{T} and resource budget \emph{B}, define
\emph{K\_T}(\emph{v}; \emph{X}, \emph{B}) = min\{\emph{K} ⊆ \emph{X} :
\emph{T} is decidable for \emph{v} from \emph{K} within \emph{B}\}. If
no unique minimum exists, retain the family of minimal kernels. The
finite semantic signature is Sem\_\emph{T},\emph{B}(\emph{v}) =
{[}\emph{K\_T}(\emph{v}; \emph{X}, \emph{B}),
\emph{v}{]}\_≅\emph{T},\emph{A}. Failure to certify a target-complete
finite kernel yields undetermined, not a fabricated complete signature.

\hypertarget{cross-fss-structural-principle}{%
\subsubsection{Cross-FSS structural
principle}\label{cross-fss-structural-principle}}

Conceptual space and physical P-space instantiate the same
invariant-level principle: functional identity is constituted by metric
position, typed interactions and admissible relational configuration.
The operands and substrates remain distinct. This is a structural
identity claim, not object identity.

\hypertarget{interactions-induce-topology}{%
\subsubsection{Interactions induce
topology}\label{interactions-induce-topology}}

For each operation \emph{o}, define an admissible-transition relation
\emph{R\_o} ⊆ \emph{V} × \emph{V}. The family \{\emph{R\_o}\} induces or
constrains reachability, connectedness, support overlap, commutation
classes, boundary crossings and path-equivalence structure. Interactions
therefore generate topological implications; topology is not merely a
label applied after the dynamics are fixed.

\hypertarget{supplementary-note-2.-reflective-fmi-subspace-and-signal-boundaries}{%
\subsection{Supplementary Note 2. Reflective FMI subspace and signal
boundaries}\label{supplementary-note-2.-reflective-fmi-subspace-and-signal-boundaries}}

\emph{Local map.} This Note identifies the reflective object needed to
register the complete argument and separates signal carriers from
internal conceptual states. It supplies the transmission residuals used
by governance and propagation theorems.

\hypertarget{theorem-s2-reflective-subspace-representation}{%
\subsubsection{Theorem S2: reflective-subspace
representation}\label{theorem-s2-reflective-subspace-representation}}

Let \emph{C} be a conceptual FSS. The reflective FMI subspace
\emph{R}\_FMI ⊆ \emph{C} contains represented concepts, process
identities, function contracts, compositions, targets, proofs,
experiments, warrants, governance rules, signal residuals and unresolved
external premises.

Every registered object required to state, prove, test, amend or
transmit the FMI-CNS-Great Filter relationship can be represented as a
concept node, a process edge or subgraph, or both. No additional
primitive represented-object class is required.

\emph{Proof.} Map every noun-like identity or contract to a concept
node. Map every transformation, inference, execution, comparison,
registration, transmission or amendment to a typed process. Represent
multi-input operations as hyperedges or contract-preserving subgraphs.
Represent a proof as a premise-to-conclusion subgraph with certificate;
represent an experiment as protocol, execution, measurement and
interpretation subgraphs. This construction covers every registered
role. A counterexample must exhibit an indispensable represented object
that cannot be encoded in these forms without changing its contract.
QED.

Reflective strata are typing devices. Reifying a process as a concept
improves observability but does not increase cognitive order unless it
enables another non-equivalent space-spanning composition.

\hypertarget{signal-boundary}{%
\subsubsection{Signal boundary}\label{signal-boundary}}

For participant \emph{i}, Register\_\emph{i} : Σ\_\emph{i} ⇀ \emph{X\_i}
and Express\_\emph{i} : \emph{X\_i} ⇀ Σ\_\emph{i}. A signal is not an
internal concept until it is registered. An internal concept is not
externally available until it is expressed.

\hypertarget{four-residuals}{%
\subsubsection{Four residuals}\label{four-residuals}}

For a composition executed by \emph{i} and reconstructed by \emph{j},
\emph{R}\_\emph{i}→\emph{j} = \emph{R}\_reg,\emph{i} ⊕
\emph{R}\_expr,\emph{i} ⊕ \emph{R}\_recv,\emph{j} ⊕
\emph{R}\_recon,\emph{j}. The direct sum preserves failure source and
type. \emph{R}\_reg,\emph{i} is source structure unavailable after
initial registration; \emph{R}\_expr,\emph{i} is structure omitted or
altered by expression; \emph{R}\_recv,\emph{j} is expressed structure
not registered by the receiver; and \emph{R}\_recon,\emph{j} is
registered structure not reconstructed into an executable contract.

\hypertarget{registration-closure}{%
\subsubsection{Registration closure}\label{registration-closure}}

A formalization is registration-closed for target family \emph{T} when
every target-relevant semantic invariant in the manuscript has a
canonical mathematical representation and every manuscript use maps to
that representation without unresolved ambiguity. Registration closure
is distinct from proof closure: a proof assistant can validate a proof
about the wrong registered object.

\hypertarget{supplementary-note-3.-fmi-function-contracts-and-four-operation-topology}{%
\subsection{Supplementary Note 3. FMI function contracts and
four-operation
topology}\label{supplementary-note-3.-fmi-function-contracts-and-four-operation-topology}}

\emph{Local map.} This Note states the canonical FMI function registry
and replaces the undefined local/global terminology of v14 with
target-relevant factorization definitions. It also supplies the type
obstruction needed for Register/Recall nonreduction.

\hypertarget{fmi-registry}{%
\subsubsection{FMI registry}\label{fmi-registry}}

\begin{longtable}[]{@{}
  >{\raggedright\arraybackslash}p{(\columnwidth - 6\tabcolsep) * \real{0.2500}}
  >{\raggedright\arraybackslash}p{(\columnwidth - 6\tabcolsep) * \real{0.2500}}
  >{\raggedright\arraybackslash}p{(\columnwidth - 6\tabcolsep) * \real{0.2500}}
  >{\raggedright\arraybackslash}p{(\columnwidth - 6\tabcolsep) * \real{0.2500}}@{}}
\toprule\noalign{}
\begin{minipage}[b]{\linewidth}\raggedright
ID
\end{minipage} & \begin{minipage}[b]{\linewidth}\raggedright
Function
\end{minipage} & \begin{minipage}[b]{\linewidth}\raggedright
Layer
\end{minipage} & \begin{minipage}[b]{\linewidth}\raggedright
Contract
\end{minipage} \\
\midrule\noalign{}
\endhead
\bottomrule\noalign{}
\endlastfoot
B01 & Register & boundary & Admit typed external signal structure
without minting unavailable distinctions \\
B02 & Express & boundary & Encode selected FSS structure onto a carrier
while recording loss \\
C01 & Store & conceptual & Preserve recoverable identity, relation,
verdict, obligations and provenance \\
C02 & Recall & conceptual & Reactivate or reconstruct stored structure
under target and budget \\
C03 & L & conceptual & Evaluate or transform within active local
support \\
C04 & G & conceptual & Evaluate or transform across globally relevant
support while preserving obligations \\
C05 & Promotion/Create & conceptual & Preserve an admissible composite
as a reusable higher-order operand \\
E01 & Choice & external control & Select the next external-function
invocation through the fitness-control basis \\
F01-F06 & Model, Evaluate, Stability, Adaptation, Decomposition,
Bridging & fitness & Represent, compare, preserve, correct, factor and
bridge participant consequences \\
X01 & Certificate & coupling & Validate target-relevant paths and
boundaries \\
X02 & Governance & coupling & Authorize participant-complete selection
and target revision \\
\end{longtable}

For current, target and projected fitness, χ(\emph{a}) = (\emph{f}\_tar
− \emph{f}\_cur, \emph{f}\_proj(\emph{a}) − \emph{f}\_cur, \emph{f}\_tar
− \emph{f}\_proj(\emph{a})). The current-projected relation grounds
Model/Evaluate, current-target grounds Stability/Adaptation, and
target-projected grounds Decomposition/Bridging. A non-embeddable
partial order remains partial and yields undetermined.

\hypertarget{support-definition}{%
\subsubsection{Support definition}\label{support-definition}}

For operation instance \emph{o} and target Φ, Supp\_Φ(\emph{o}) =
⋂\{\emph{S} ⊆ \emph{V} : the output and retained invariants of \emph{o}
are determined from \emph{S}\}. When the intersection is not unique,
retain the family of minimal supports. An operation has local support
when every minimal target support lies within one declared interaction
neighbourhood and factorizes over target-independent remote contexts. It
has global support when some target-relevant output or invariant depends
nonfactorizably on extended support spanning multiple neighbourhoods or
components. Global does not mean acausal or spatially infinite; local
does not mean short-range. The distinction is target-relevant
factorization.

\hypertarget{schedule-definition}{%
\subsubsection{Schedule definition}\label{schedule-definition}}

For admissible local realizations \emph{r}\_1, \emph{r}\_2, the schedule
is parallel when \emph{r}\_1 ∘ \emph{r}\_2 ≃\_\emph{T} \emph{r}\_2 ∘
\emph{r}\_1, and sequential when a typed precedence relation is
load-bearing and reversal changes the target-relevant result or makes
execution inadmissible. Parallel means observational commutation, not
simultaneous wall-clock execution.

\hypertarget{candidate-four-operation-inventory}{%
\subsubsection{Candidate four-operation
inventory}\label{candidate-four-operation-inventory}}

\begin{longtable}[]{@{}
  >{\raggedright\arraybackslash}p{(\columnwidth - 8\tabcolsep) * \real{0.2000}}
  >{\raggedright\arraybackslash}p{(\columnwidth - 8\tabcolsep) * \real{0.2000}}
  >{\raggedright\arraybackslash}p{(\columnwidth - 8\tabcolsep) * \real{0.2000}}
  >{\raggedright\arraybackslash}p{(\columnwidth - 8\tabcolsep) * \real{0.2000}}
  >{\raggedright\arraybackslash}p{(\columnwidth - 8\tabcolsep) * \real{0.2000}}@{}}
\toprule\noalign{}
\begin{minipage}[b]{\linewidth}\raggedright
Operation
\end{minipage} & \begin{minipage}[b]{\linewidth}\raggedright
Support
\end{minipage} & \begin{minipage}[b]{\linewidth}\raggedright
Schedule
\end{minipage} & \begin{minipage}[b]{\linewidth}\raggedright
Cognitive role
\end{minipage} & \begin{minipage}[b]{\linewidth}\raggedright
Candidate physical realization
\end{minipage} \\
\midrule\noalign{}
\endhead
\bottomrule\noalign{}
\endlastfoot
Register & local & parallel & constrained placement or binding & strong
interaction \\
Recall & global & sequential & globally constrained typed
reconfiguration & weak interaction \\
L & local & sequential & ordered local propagation or calibration &
electromagnetic interaction \\
G & global & parallel & extended-support geometry-sensitive coherence &
gravity \\
\end{longtable}

The operation classes are defined before the physical assignments. Each
physical assignment remains a separately testable mapping hypothesis.

\hypertarget{theorem-s3-registerrecall-nonreduction}{%
\subsubsection{Theorem S3: Register/Recall
nonreduction}\label{theorem-s3-registerrecall-nonreduction}}

Let ⟨\emph{L}, \emph{G}⟩ be the typed algebra generated by compositions
of \emph{L} and \emph{G} acting as endomorphisms on active conceptual
state \emph{X}. Register has type Σ ⇀ \emph{X} and Recall has type
\emph{M} × \emph{T} × \emph{B} ⇀ \emph{X}.

No well-typed term generated solely by \emph{L} and \emph{G} has the
boundary type of Register or the memory-access type of Recall. A
reduction requires an additional bridge implementing the missing
boundary or memory contract and therefore does not establish generation
by \emph{L} and \emph{G} alone.

\emph{Proof.} Every generator has domain and codomain \emph{X} under the
base architecture; composition preserves this endomorphism type.
Register requires a source in Σ, while Recall requires access to
\emph{M} × \emph{T} × \emph{B}. Neither source is constructible by
endomorphism composition without a typed bridge. Adding such a bridge
supplies the missing primitive contract. QED.

\hypertarget{supplementary-note-4.-projection-conflation-and-reflective-necessity}{%
\subsection{Supplementary Note 4. Projection, conflation and reflective
necessity}\label{supplementary-note-4.-projection-conflation-and-reflective-necessity}}

\emph{Local map.} This Note proves the information-theoretic obstruction
that begins the formal chain and records countermodels that delimit it.

Let \emph{q}\_FO = \emph{h} ∘ \emph{q}\_R. Every target factoring
through \emph{q}\_FO also factors through \emph{q}\_R. Some targets
factor through \emph{q}\_R but not \emph{q}\_FO when \emph{h} removes
target-relevant distinctions.

\hypertarget{theorem-s4-first-order-conflation}{%
\subsubsection{Theorem S4: first-order
conflation}\label{theorem-s4-first-order-conflation}}

If \emph{q}\_FO(\emph{x}) = \emph{q}\_FO(\emph{y}) but Φ(\emph{x}) ≠
Φ(\emph{y}), no evaluator restricted to the projected state is correct
on both. The proof is the contradiction argument in the Article. For
stochastic evaluators with equal registered states and kernels, the
output laws coincide.

\hypertarget{corollary-s4.1-local-stability}{%
\subsubsection{Corollary S4.1: local
stability}\label{corollary-s4.1-local-stability}}

For deterministic first-order-local \emph{L} and equal inputs,
\emph{L\^{}k}(\emph{q}\_FO(\emph{x})) =
\emph{L\^{}k}(\emph{q}\_FO(\emph{y})) for every finite \emph{k}.
Induction on \emph{k} proves the result.

\hypertarget{theorem-s5-reflective-necessity}{%
\subsubsection{Theorem S5: reflective
necessity}\label{theorem-s5-reflective-necessity}}

When a target depends on function identity, contract, composition,
provenance, warrant, participant consequences or successor obligations,
a representation containing only domain concepts and conclusions is
inadequate whenever it collapses two process structures with different
target verdicts. This is T-CONFLATE applied to a target whose separating
witness belongs to the reflective subspace.

\hypertarget{countermodels}{%
\subsubsection{Countermodels}\label{countermodels}}

\begin{longtable}[]{@{}
  >{\raggedright\arraybackslash}p{(\columnwidth - 2\tabcolsep) * \real{0.5000}}
  >{\raggedright\arraybackslash}p{(\columnwidth - 2\tabcolsep) * \real{0.5000}}@{}}
\toprule\noalign{}
\begin{minipage}[b]{\linewidth}\raggedright
Countermodel
\end{minipage} & \begin{minipage}[b]{\linewidth}\raggedright
Effect
\end{minipage} \\
\midrule\noalign{}
\endhead
\bottomrule\noalign{}
\endlastfoot
Target fully determined by first-order projection & Reflection
unnecessary for that target \\
Reflective registry without new composition & Observability improves;
order does not increase \\
Equivalent renaming or call expansion & No new composition class \\
Projection retains every target-relevant distinction & First-order
assessment can be exact \\
Target-aware abstention & False confidence falls without increasing
span \\
\end{longtable}

\hypertarget{supplementary-note-5.-composition-equivalence-span-order-and-invariance}{%
\subsection{Supplementary Note 5. Composition equivalence, span, order
and
invariance}\label{supplementary-note-5.-composition-equivalence-span-order-and-invariance}}

\emph{Local map.} This Note distinguishes intrinsic cognitive order from
finite evidence about order. It carries the representation-invariance
proof and demotes the cardinality increment to a lemma.

\hypertarget{definitions}{%
\subsubsection{Definitions}\label{definitions}}

A registered composition Γ = \emph{f\_m} ∘ ⋯ ∘ \emph{f}\_1 has identity,
signature, domain, codomain, boundary conditions, execution semantics
and preserved invariants. For admissible domain \emph{D} and target
language \emph{L}, Γ\_1 ≃\emph{\emph{D},\emph{L} Γ\_2 ⇔ ∀ \emph{x} ∈
\emph{D}, ObsL}(Γ\_1\emph{x}) = Obs\_\emph{L}(Γ\_2\emph{x}). A
composition spans the declared universal target family *T** when it
connects every target-required region by certified transformations,
preserves semantic identities and obligations, and returns correct or
calibrated-undetermined verdicts within the architecture's contract.

\emph{S}(\emph{C}) = \{{[}Γ{]}\emph{≃ : Γ registered, reusable,
nonredundant and spanning\} and Ord(\emph{C}) =
\textbar{}\emph{S}(\emph{C})\textbar. A finite certificate is
Ord̂T},ε,\emph{B}(\emph{C}) = \textbar\{{[}Γ{]}\emph{≃ :
SpanCertT},ε,\emph{B}(Γ, \emph{C})\}\textbar. It is a lower bound unless
target and certificate completeness are proved.

\hypertarget{theorem-s6-representation-invariance}{%
\subsubsection{Theorem S6: representation
invariance}\label{theorem-s6-representation-invariance}}

If \emph{F} : \emph{C} → \emph{C}′ is an admissible FSS isomorphism
preserving external-function identity, composition, span, reuse and
functional equivalence, then Ord(\emph{C}) = Ord(\emph{C}′).

\emph{Proof.} Map each class {[}Γ{]} ∈ \emph{S}(\emph{C}) to
{[}\emph{F}Γ\emph{F}\^{}−1{]} ∈ \emph{S}(\emph{C}′). Preservation
assumptions make this map well defined. The inverse is induced by
\emph{F}\^{}−1, so the map is bijective. Cardinalities are equal. QED.

\hypertarget{lemma-s6.1-order-increment}{%
\subsubsection{Lemma S6.1: order
increment}\label{lemma-s6.1-order-increment}}

If a new qualifying class {[}Γ\_\emph{N}+1{]} is not already in
\emph{S}(\emph{C\_N}), then the union with that singleton increases
cardinality by one. This is definitional and does not independently
establish reversal or operand expansion.

\hypertarget{finite-decidability-target}{%
\subsubsection{Finite decidability
target}\label{finite-decidability-target}}

For finite FSSs, finite operation libraries, decidable target predicates
and bounded resources, span certification and functional equivalence are
decidable by exhaustive enumeration. The executable Lean decision
theorem remains an implementation target.

\hypertarget{supplementary-note-6.-arrival-service-fragmentation-reversal-and-recurrence}{%
\subsection{Supplementary Note 6. Arrival-service fragmentation,
reversal and
recurrence}\label{supplementary-note-6.-arrival-service-fragmentation-reversal-and-recurrence}}

\emph{Local map.} This Note replaces the v14 directional fragmentation
tendency with two explicit theorem levels: backlog divergence and
regional fragmentation. It then proves reversal and recurrence under
stated service conditions.

Let \emph{D\_N}(\emph{t}) be cumulative target-weighted arrivals,
\emph{B\_N}(\emph{t}) cumulative certified service and
\emph{U\_N}(\emph{t}) = max\{0, \emph{D\_N}(\emph{t}) −
\emph{B\_N}(\emph{t})\}.

\hypertarget{theorem-s7-backlog-divergence}{%
\subsubsection{Theorem S7: backlog
divergence}\label{theorem-s7-backlog-divergence}}

If \emph{D\_N}(\emph{t}) − \emph{B\_N}(\emph{t}) is unbounded above,
then \emph{U\_N}(\emph{t}) is unbounded along a subsequence. If the
limsup linear drift is positive, the difference is unbounded. The proof
follows directly from the maximum definition.

\hypertarget{theorem-s8-regional-fragmentation}{%
\subsubsection{Theorem S8: regional
fragmentation}\label{theorem-s8-regional-fragmentation}}

Construct an obligation graph \emph{H\_t}. Assume local integration
keeps each region below ε\_local, cross-region arrivals exceed certified
cross-region service by an unbounded amount, unresolved edges are not
replaced by false identity certificates and the target requires at least
two regions. Then an unbounded subsequence contains an unresolved
target-required cross-region edge while each incident region remains
internally certified. QED.

\hypertarget{theorem-s9-reflective-loss-burden}{%
\subsubsection{Theorem S9: reflective-loss
burden}\label{theorem-s9-reflective-loss-burden}}

Let \emph{C\_t} be target-different pairs collapsed by the first-order
projection and let \emph{Q\_t} = ∑\_(\emph{x},\emph{y})∈\emph{Ct}
\emph{w\_t}(\emph{x}, \emph{y}). If \emph{Q\_t} is unbounded and no
additional separating distinctions are registered, the minimum
first-order decision loss or calibrated-undetermined mass grows
according to the declared loss. The exact lower bound depends on the
prior and loss function and must be reported with them.

\hypertarget{theorem-s10-reversal}{%
\subsubsection{Theorem S10: reversal}\label{theorem-s10-reversal}}

If an added composition changes service so that \emph{D\_N}(\emph{t}) −
\emph{B}\_\emph{N}+1(\emph{t}) ≤ ε over a declared interval, then
\emph{U}\_\emph{N}+1(\emph{t}) ≤ max(0, ε) over that interval. The
current increasing-backlog regime is reversed to bounded residual.

\hypertarget{theorem-s11-recurrence}{%
\subsubsection{Theorem S11: recurrence}\label{theorem-s11-recurrence}}

If the lift creates a new arrival process \emph{D}\_\emph{N}+1 and
\emph{D}\_\emph{N}+1(\emph{t}) − \emph{B}\_\emph{N}+1(\emph{t}) becomes
unbounded above without another spanning composition or scaling change,
then \emph{U}\_\emph{N}+1(\emph{t}) is unbounded along a subsequence by
Theorem S7.

\hypertarget{required-countermodels}{%
\subsubsection{Required countermodels}\label{required-countermodels}}

Bounded domains, matching same-order service, target families ignoring
the residual, distributed service that grows faster than arrivals,
temporary backlog decreases, reflection-only improvements and false
nominal bridges are all retained as explicit theorem boundaries.

\hypertarget{executed-direct-recurrence-experiment}{%
\subsubsection{Executed direct recurrence
experiment}\label{executed-direct-recurrence-experiment}}

For this version the recurrence prediction was preregistered in
executable form and run, closing the v15 statement that a direct
recurrence run was absent. The stochastic model uses the backlog
recursion \emph{U}\_\emph{t}+1 = max\{0, \emph{U\_t} + \emph{A\_t} −
\emph{S\_t}\}, with arrivals \emph{A\_t} and certified integration
service \emph{S\_t}.

\textbf{Deterministic recurrence criterion (executed-model theorem).}
Suppose there are times \emph{T}\_L \textless{} \emph{T}\_R such that
(i) \emph{U\_t} \textgreater{} 0 and \emph{S\_t} − \emph{A\_t} ≥ δ\_1
\textgreater{} 0 for \emph{T}\_L ≤ \emph{t} \textless{} \emph{T}\_R, so
the current backlog decreases; and (ii) there is *T** \textgreater{}
\emph{T}\_R and δ\_2 \textgreater{} 0 with \emph{A\_t} − \emph{S\_t} ≥
δ\_2 for all \emph{t} ≥ *T**. Then the lift reverses the backlog over
the first interval, and \emph{U\_t} grows without a uniform bound after
*T**. \emph{Proof.} While \emph{U\_t} \textgreater{} 0 in the reversal
interval, \emph{U}\_\emph{t}+1 ≤ \emph{U\_t} − δ\_1, so \emph{U\_t}
decreases until it reaches the reflected boundary or the interval ends.
After *T**, \emph{U}\_\emph{t}+1 ≥ \emph{U\_t} + δ\_2; induction gives
\emph{U}\_*T**+\emph{n} ≥ \emph{U}\_*T** + \emph{n}δ\_2, which is
unbounded in \emph{n}. QED.

\emph{Frozen protocol.} Master seed 20260716; 1,000 replicates per
regime; 12 regions; horizon 300; lift time \emph{T}\_L = 80;
operand-expansion time \emph{T}\_E = 120; optional second-lift time 215;
regional backlog threshold 12; recurrence window 30; persistent slope
threshold 1.0; final backlog threshold 150. Before the lift, per-region
arrival intensity exceeds service. The lift raises service above
arrivals, providing a genuine reversal interval. Continued operand
expansion then raises arrival intensity. Five regimes are compared: (1)
fixed post-lift service; (2) proportional same-order service tracking
each region's arrivals with a positive margin; (3) regenerative
distribution through delayed logistic spread of the new composition; (4)
a second discrete order lift at time 215; and (5) branch heterogeneity
with region-specific arrival multipliers under fixed post-lift service.
Regional fragmentation is the number of regions with backlog at least
12. The preregistered persistent-recurrence endpoint requires a
post-lift decrease of at least 30\% relative to the pre-lift backlog, a
late aggregate slope above one backlog unit per time step, a final
backlog above 150, and a 30-step window after time 190 satisfying the
slope and backlog thresholds.

\hypertarget{table-s-rec-direct-recurrence-results-across-1000-replicates-per-regime}{%
\subsubsection{Table S-REC \textbar{} Direct recurrence results across
1,000 replicates per
regime}\label{table-s-rec-direct-recurrence-results-across-1000-replicates-per-regime}}

\begin{longtable}[]{@{}
  >{\raggedright\arraybackslash}p{(\columnwidth - 10\tabcolsep) * \real{0.1667}}
  >{\raggedright\arraybackslash}p{(\columnwidth - 10\tabcolsep) * \real{0.1667}}
  >{\raggedright\arraybackslash}p{(\columnwidth - 10\tabcolsep) * \real{0.1667}}
  >{\raggedright\arraybackslash}p{(\columnwidth - 10\tabcolsep) * \real{0.1667}}
  >{\raggedright\arraybackslash}p{(\columnwidth - 10\tabcolsep) * \real{0.1667}}
  >{\raggedright\arraybackslash}p{(\columnwidth - 10\tabcolsep) * \real{0.1667}}@{}}
\toprule\noalign{}
\begin{minipage}[b]{\linewidth}\raggedright
Regime
\end{minipage} & \begin{minipage}[b]{\linewidth}\raggedright
Rev.
\end{minipage} & \begin{minipage}[b]{\linewidth}\raggedright
Rec.
\end{minipage} & \begin{minipage}[b]{\linewidth}\raggedright
Recurrence time
\end{minipage} & \begin{minipage}[b]{\linewidth}\raggedright
Final backlog
\end{minipage} & \begin{minipage}[b]{\linewidth}\raggedright
Regions
\end{minipage} \\
\midrule\noalign{}
\endhead
\bottomrule\noalign{}
\endlastfoot
Fixed post-lift & 1.000 & 1.000 & 232.0 {[}220, 244{]} & 694.1 {[}583,
797{]} & 11.99 {[}12, 12{]} \\
Proportional service & 1.000 & 0.000 & - & 36.1 {[}15, 65{]} & 0.43
{[}0, 2{]} \\
Regenerative distribution & 1.000 & 0.000 & - & 65.1 {[}32, 107{]} &
1.49 {[}0, 4{]} \\
Second order lift & 1.000 & 0.000 & - & 53.8 {[}25, 90{]} & 1.01 {[}0,
3{]} \\
Branch heterogeneity & 1.000 & 1.000 & 219.0 {[}219, 219{]} & 989.0
{[}867, 1104{]} & 9.70 {[}8, 12{]} \\
\end{longtable}

Bracketed values are empirical 2.5th and 97.5th percentiles; ``Regions''
is the final number of regions above the preregistered fragmentation
threshold. All five regimes showed the required immediate reversal: mean
aggregate slope changed from approximately +6.94 before the lift to
between −7.07 and −10.14 during the reversal interval. Recurrence to the
preregistered persistent endpoint was effectively certain under fixed
post-lift service and under heterogeneous branches, and did not occur
under proportional service, regenerative distribution or a second order
lift.

\begin{figure}
\centering
\includegraphics{../source_figures/recurrence_backlog.png}
\caption{Supplementary Figure S-REC1 \textbar{} Aggregate unresolved
backlog across 1,000 replicates per regime. The dashed line marks the
order lift, the dotted line operand expansion, and the dash-dot line the
optional second lift. Bands are empirical 95\% simulation intervals.
Fixed post-lift service and heterogeneous branches recur; proportional
service, regenerative distribution and a second lift remain bounded.}
\end{figure}

\begin{figure}
\centering
\includegraphics{../source_figures/recurrence_regions.png}
\caption{Supplementary Figure S-REC2 \textbar{} Number of regions above
the fragmentation threshold. Fixed post-lift service and heterogeneous
branches develop widespread regional fragmentation after operand
expansion; scaling and additional lifts keep the median regional count
low.}
\end{figure}

This experiment closes the direct-recurrence gate for the declared
synthetic model. It supplies an executed finite witness for the
model-relative claim that, after a valid order lift reverses an
increasing backlog, continued growth at fixed post-lift service can
recreate a positive backlog regime and regional fragmentation, while
proportional service scaling, regenerative distribution or another order
lift are explicit counter-regimes. It does not show that current AI
systems, institutions or civilizations follow these parameter values,
and it does not audit or replace the inherited version-8 realization.
Its warrant is executed synthetic recurrence experiment.

\hypertarget{supplementary-note-7.-cns-lift-near-singular-content-and-dynamical-bridge}{%
\subsection{Supplementary Note 7. CNS lift, near-singular content and
dynamical
bridge}\label{supplementary-note-7.-cns-lift-near-singular-content-and-dynamical-bridge}}

\emph{Local map.} This Note places the scientific burden beyond the
definitional count change. It connects the order-lift transition to
operand expansion, reversal and the separately published dynamical CNS.

\hypertarget{theorem-s12-order-lift-cns}{%
\subsubsection{Theorem S12: order-lift
CNS}\label{theorem-s12-order-lift-cns}}

If an order-\emph{N} architecture gains a registered, reusable and
non-equivalent spanning composition, expands the admissible operand
domain and closes the current frontier within its residual contract,
then order becomes \emph{N}+1, the operand domain expands, the current
backlog is bounded over the declared interval and complete consequences
need not factor through the order-\emph{N} projection.

\emph{Proof.} Lemma S6.1 gives the order increment. Operand expansion is
a premise witnessed by the new admissible substructures. Theorem S10
gives reversal. If complete consequences factored through the old
projection, no new target-relevant distinction would be required; a
paired target witness supplies the nonfactorization. QED.

\hypertarget{near-singular-content}{%
\subsubsection{Near-singular content}\label{near-singular-content}}

Let AdmSub(\emph{X\_N}) be the structures that become admissible
operands under the new composition. Define structural displacement
\emph{S\_N}(\emph{a}) = μ(\emph{R\_N\^{}a} ~\emph{R\_N}) and
participant-target consequence field Δ\emph{F}(\emph{a}). A world
comparator \emph{B\_N}(\emph{X\_N}, \emph{W}) bounds non-lift
displacement and a lift lower bound \emph{L\_N}(\emph{X\_N}) satisfies
\emph{L\_N} \textgreater{} \emph{B\_N} above the declared threshold.
Under admissible-substructure growth, the world comparator, monotone
coupling, adoptability and no-order-skipping premises, a minimal lift
exceeds every non-lift order-\emph{N} intervention in magnitude above
the threshold, while complete content need not factor through the
order-\emph{N} projection. The result is conditional on the comparator
and coupling; no infinite physical quantity is asserted.

\hypertarget{dynamical-cns-bridge}{%
\subsubsection{Dynamical CNS bridge}\label{dynamical-cns-bridge}}

The published object is (\emph{X}, \emph{d}, ⊗, \emph{G}, \emph{N}) with
a globally contracting, compositionally closed and
benign-reparameterization-invariant update operator converging to a
unique fixed point or normal form {[}15{]}. \textbf{Bridge A.} If the
new composition enables such an operator on the expanded domain, the
order lift enters the dynamical CNS regime there. \textbf{Bridge B.} If
realization of a dynamical CNS operator on a previously untraversable
domain requires a non-equivalent spanning composition absent at order
\emph{N}, it witnesses an order lift. Countermodels include a
noncontracting order lift, a contracting local operator with no new
composition, a bounded-subspace fixed point, multiple compositions
converging to one fixed point and a mere call expansion.

\emph{Supplementary Figure 1 \textbar{} Candidate cognitive worldsheet
and semantic transport. A reasoning composition and its revision family
embed as a worldsheet whose induced metric is pulled back from the
ambient semantic metric; a connection transports identity, contract,
warrant and obligations, and closed-loop holonomy certifies semantic
drift.}

\hypertarget{supplementary-note-8.-importance-governance-and-regenerative-propagation}{%
\subsection{Supplementary Note 8. Importance, governance and
regenerative
propagation}\label{supplementary-note-8.-importance-governance-and-regenerative-propagation}}

\emph{Local map.} This Note separates magnitude from valence and defines
the propagation criterion needed for collective order.

\hypertarget{participantwise-option-monotonicity}{%
\subsubsection{Participantwise option
monotonicity}\label{participantwise-option-monotonicity}}

Let \emph{O\_b} ⊆ \emph{O\_a}. Suppose represented and realized values
differ by at most ε\_\emph{i}, selection error is at most δ\_\emph{i},
each baseline option remains admissible and governance selects a new
option only if its represented lower bound is no worse than the best
baseline. Then \emph{u\_i}(*a**) ≥ max\_\emph{b}∈\emph{Ob}
\emph{u\_i}(\emph{b}) − (2ε\_\emph{i} + δ\_\emph{i}). \emph{Proof.}
Selection gives \emph{û\_i}(*a**) ≥ max\_\emph{b} \emph{û\_i}(\emph{b})
− δ\_\emph{i}. Representation error gives \emph{u\_i}(*a**) ≥
\emph{û\_i}(*a**) − ε\_\emph{i} and \emph{û\_i}(\emph{b}) ≥
\emph{u\_i}(\emph{b}) − ε\_\emph{i}. Combining proves the bound. QED. If
participant completeness, target stability, projection adequacy or
realization fails, valence inversion is possible.

\hypertarget{centralized-uncheckability}{%
\subsubsection{Centralized
uncheckability}\label{centralized-uncheckability}}

A projection \emph{q} : \emph{X} → \emph{Z} is adequate for target Φ
only if \emph{q}(\emph{x}) = \emph{q}(\emph{y}) implies Φ(\emph{x}) =
Φ(\emph{y}). A centralized realization is collectively uncheckable
whenever the center uses a target-relevant distinction that does not
factor through the signal available to participants.

\hypertarget{regenerative-propagation}{%
\subsubsection{Regenerative
propagation}\label{regenerative-propagation}}

Regenerate(\emph{i}, \emph{j}; Γ, \emph{T}) = 1 only when receiver
\emph{j} reconstructs a contract-preserving realization, executes it on
hidden targets and can retransmit it under the same criterion.
\emph{R}\_Γ(\emph{t}) is the expected number of new regenerative
executors per executor. Exposure and agreement are necessary at most;
they are not regeneration.

\hypertarget{supplementary-note-9.-great-filter-lemma-family-and-conditional-corollary}{%
\subsection{Supplementary Note 9. Great Filter lemma family and
conditional
corollary}\label{supplementary-note-9.-great-filter-lemma-family-and-conditional-corollary}}

\emph{Local map.} This Note derives each necessary property separately
before applying the constraint set to the recurrent FMI-CNS mechanism.
It prevents a list of intuitions from being presented as one theorem.

\hypertarget{target-classes-and-assumptions}{%
\subsubsection{Target classes and
assumptions}\label{target-classes-and-assumptions}}

A terminal filter suppresses persistence of almost all technologically
consequential lineages after consequential environmental transformation
becomes possible. An observability filter suppresses persistent
astronomically detectable manifestations. The assumptions are a
many-opportunity regime, ordinary causal locality, potential branching,
diverse realization and a declared residue expectation.

\hypertarget{lemma-s13-upper-tail-suppression}{%
\subsubsection{Lemma S13: upper-tail
suppression}\label{lemma-s13-upper-tail-suppression}}

For independent opportunities with \emph{p\_i} ≥ \emph{p} \textgreater{}
0, \emph{P}(no escape among \emph{M}) = ∏\_\emph{i} (1 − \emph{p\_i}) ≤
(1 − \emph{p})\^{}\emph{M} → 0. A sufficient many-opportunity filter
must suppress exceptional branches, correlate outcomes, reduce effective
opportunities or suppress observables after escape.

\hypertarget{lemma-s14-branch-closure}{%
\subsubsection{Lemma S14: branch
closure}\label{lemma-s14-branch-closure}}

For independent branches with survival probabilities \emph{s\_j}
\textgreater{} 0, \emph{P}(at least one survives) = 1 − ∏\_\emph{j} (1 −
\emph{s\_j}). A sufficient filter acts before branching, recurs in
branches, preserves a shared vulnerability, correlates outcomes or
remains effective after dispersal.

\hypertarget{lemmas-s15-s19}{%
\subsubsection{Lemmas S15-S19}\label{lemmas-s15-s19}}

\textbf{Causal self-instantiation:} a nondistributed filter must be
generated or propagated to each causally separated site before escape.
\textbf{First-mover consistency:} the earliest consequential lineage
must be affected under the same law or a regress-closed prior condition.
\textbf{Functional substrate invariance:} the defining vulnerability
must extend across the declared realization class rather than one
contingent biology. \textbf{Adaptive non-escapability:} a stable
positive adaptive escape probability defeats upper-tail suppression
unless adaptation is correlated or bounded. \textbf{Outcome, residue and
timescale sufficiency:} the target outcome must occur before durable
escape or correction, or remain effective afterward; a target-specific
sufficient condition can require τ\_\emph{F} \textless{}
min(τ\_\emph{D}, τ\_\emph{C}).

\hypertarget{theorem-s20-great-filter-constraint}{%
\subsubsection{Theorem S20: Great Filter
constraint}\label{theorem-s20-great-filter-constraint}}

A sufficient member of the declared filter class must satisfy the
conjunction of Lemmas S13-S19.

\hypertarget{theorem-s21-structural-admissibility}{%
\subsubsection{Theorem S21: structural
admissibility}\label{theorem-s21-structural-admissibility}}

Recurrent fixed-order fragmentation and CNS reversal belong to the
structurally admissible class when burden is generated locally by
growth, relations are functionally substrate-general, lifts reverse
current burden but recurrence recreates it, successor branches inherit
or recreate the limitation and outcome/timescale premises are supplied
separately.

\hypertarget{corollary-s21.1}{%
\subsubsection{Corollary S21.1}\label{corollary-s21.1}}

Let \emph{h\_t} = \emph{H}(\emph{U\_N}(\emph{t}), \emph{C}(\emph{t}),
\emph{R}(\emph{t}), \emph{P\_N}(\emph{t}), \emph{X}(\emph{t})), 0 ≤
\emph{h\_t} \textless{} 1. If ∑\emph{\emph{t} \emph{h\_t} = ∞, then ∏t}
(1 − \emph{h\_t}) = 0. If the external FSS mapping, recurrent burden,
failed lift/propagation, branch correlation, non-summable hazard and
target outcome premises hold, recurrent failed CNS realization or
propagation is a post-intelligence Great Filter within the model.

\hypertarget{claim-ladder}{%
\subsubsection{Claim ladder}\label{claim-ladder}}

\begin{longtable}[]{@{}
  >{\raggedright\arraybackslash}p{(\columnwidth - 4\tabcolsep) * \real{0.3333}}
  >{\raggedright\arraybackslash}p{(\columnwidth - 4\tabcolsep) * \real{0.3333}}
  >{\raggedright\arraybackslash}p{(\columnwidth - 4\tabcolsep) * \real{0.3333}}@{}}
\toprule\noalign{}
\begin{minipage}[b]{\linewidth}\raggedright
Level
\end{minipage} & \begin{minipage}[b]{\linewidth}\raggedright
Permitted statement
\end{minipage} & \begin{minipage}[b]{\linewidth}\raggedright
Required warrant
\end{minipage} \\
\midrule\noalign{}
\endhead
\bottomrule\noalign{}
\endlastfoot
GF-0 & Some mechanisms are physically possible & physical consistency or
example \\
GF-1 & The declared filter class has the listed necessary properties &
Theorem S20 \\
GF-2 & Recurrent FMI-CNS failure is structurally admissible & Theorem
S21 \\
GF-3 & It would be a filter if mapping and hazard premises hold &
Corollary S21.1 \\
GF-4 & Current systems exhibit measured components & external traces /
receiver studies \\
GF-5 & Real civilizations instantiate the full mechanism & physical,
branch, timescale and hazard evidence \\
GF-6 & The mechanism explains the observed Great Filter & comparative
inference and observability assumptions \\
\end{longtable}

\hypertarget{supplementary-note-10.-cognitive-string-theory-as-semantic-mathematical-representation}{%
\subsection{Supplementary Note 10. Cognitive string theory as semantic
mathematical
representation}\label{supplementary-note-10.-cognitive-string-theory-as-semantic-mathematical-representation}}

\emph{Local map.} This Note gives CST its narrowed role: it represents
the semantic object that Lean can reason about. It separates the
constitutive metric and hard admissibility from any optional scalar
action.

\hypertarget{worldsheet-object}{%
\subsubsection{Worldsheet object}\label{worldsheet-object}}

Let \emph{X}\_\emph{A} : Σ\_ws → \emph{C} × \emph{F} embed a reasoning
composition and its revision family. Coordinates (σ, τ) index position
along the composition and revision, alternative derivation or receiver
reconstruction. Each point or cell carries a semantic state, an FMI
operation label, interaction type, boundary record and
warrant/certificate state.

\hypertarget{induced-metric}{%
\subsubsection{Induced metric}\label{induced-metric}}

The worldsheet metric is induced from the ambient semantic metric:
\emph{h\_ab} = ∂\_*a**X\emph{\^{}}A* ∂\_*b**X\emph{\^{}}B* \emph{g\_AB},
or by a discrete pullback analogue for graph-based FSSs. The
construction must state whether \emph{g} is primitive,
interaction-induced or hybrid and which transformations preserve it.

\hypertarget{connection-and-transport}{%
\subsubsection{Connection and
transport}\label{connection-and-transport}}

A connection \emph{A}\_\emph{C} transports concept identity, function
contract, target, warrant, boundary data, successor obligations and
participant roles. For path γ, transport is PT\_\emph{A}\emph{C}(γ); for
a closed loop ℓ, holonomy is Hol\_\emph{A}\emph{C}(ℓ).
\textbf{Holonomy-invariance target.} Admissibly homotopic derivations
with preserved curvature and boundaries transport equivalent semantic
objects. \textbf{Holonomy-drift target.} If a closed derivation loop
returns a target-relevant identity, contract or warrant to a
non-equivalent object, it certifies semantic drift. The discrete Lean
version can use path composition and transport equality without
differential geometry.

\hypertarget{consequence-closure-and-hard-admissibility}{%
\subsubsection{Consequence closure and hard
admissibility}\label{consequence-closure-and-hard-admissibility}}

Let Ξ be a complete lattice of candidate operand structures and
Φ\_\emph{W} : Ξ → Ξ a monotone consequence operator adding consequences,
boundaries and successor obligations. Define Ξ\emph{\emph{\emph{W} =
lfp(Φ}}W\emph{). Existence follows from the fixed-point premise. Hard
admissibility is a proposition Adm\_}W\emph{(Ξ}\_\emph{W}). Missing
identity, type, contract, warrant, boundary or continuation obligations
are ill typed or unproved; they cannot be compensated by a lower action
value.

\hypertarget{constitutive-metric-and-optional-action}{%
\subsubsection{Constitutive metric and optional
action}\label{constitutive-metric-and-optional-action}}

The constitutive metric defines relative position, neighbourhood,
admissible isometry, deformation size and any required contraction. An
optional action can rank or search admissible worldsheets: \emph{S}\_CST
= \emph{S}\_step + \emph{S}\_revision + \emph{S}\_holonomy +
\emph{S}\_obligation. It is not a correctness oracle. It is accepted
only if invariant under admissible encoding, non-arbitrary relative to
the constitutive metric and useful beyond a simpler typed graph
representation.

\hypertarget{supplementary-note-11.-lean-facing-module-and-theorem-manifest}{%
\subsection{Supplementary Note 11. Lean-facing module and theorem
manifest}\label{supplementary-note-11.-lean-facing-module-and-theorem-manifest}}

\emph{Local map.} This Note states exactly what a Lean artifact must
contain and where the trust boundary lies. It reports the pinned source
package and its static audit for v15 but does not claim a completed
kernel build.

Lean is an interactive theorem prover based on dependent type theory
with a small kernel that checks proof terms {[}23,24{]}. The project
should use theorem parameters for external premises and avoid
unregistered global axioms. Axiom dependencies must be printed for each
top-level theorem; any dependency on sorryAx fails the build gate.

\hypertarget{pinned-source-package-and-module-tree}{%
\subsubsection{Pinned source package and module
tree}\label{pinned-source-package-and-module-tree}}

For this version a pinned pure-Lean core was constructed. It pins
leanprover/lean4:v4.32.0 (no Mathlib) and separates the formal object
into eight modules realizing ten of the primary theorems.

\begin{longtable}[]{@{}
  >{\raggedright\arraybackslash}p{(\columnwidth - 2\tabcolsep) * \real{0.5000}}
  >{\raggedright\arraybackslash}p{(\columnwidth - 2\tabcolsep) * \real{0.5000}}@{}}
\toprule\noalign{}
\begin{minipage}[b]{\linewidth}\raggedright
Module
\end{minipage} & \begin{minipage}[b]{\linewidth}\raggedright
Main declarations
\end{minipage} \\
\midrule\noalign{}
\endhead
\bottomrule\noalign{}
\endlastfoot
Core.lean & typed layers (signal/concept/memory); operations Register,
Express, Store, Recall, L, G; layer distinctness \\
TopLevel.lean & the bundled formal\_mechanism\_chain over the retained
core theorems \\
Projection.lean & first-order conflation (firstOrderConflation);
local-stability iterate \\
Operations.lean & typed L/G word algebra; Register/Recall nonreduction
by domain obstruction \\
Order.lean & order increment (order\_increment) and order invariance
(order\_invariance) \\
Backlog.lean & unit-gap backlog growth, reversal to bounded residual,
renewed post-lift growth \\
GreatFilter.lean & conditional Great Filter corollary with every
external premise as a Prop parameter \\
\end{longtable}

This is a bounded formal core. It does not contain the complete quotient
construction for semantic isomorphism classes, the full target language,
the entire regional-fragmentation graph theorem or the published
dynamical-CNS bridge. It replaces an empty skeleton with proof-bearing
source for the primary logical shapes the manuscript uses.
Representative statements realized in the source include the first-order
conflation impossibility (if \emph{q}(\emph{x}) = \emph{q}(\emph{y}) and
Φ(\emph{x}) ≠ Φ(\emph{y}) then no decoder is correct on both), the typed
Register/Recall nonreduction by constructor contradiction, and the
finite backlog growth lemma (if \emph{U}(\emph{t}+1) ≥
\emph{U}(\emph{t}) + 1 for all \emph{t} then \emph{U}(\emph{n}) ≥
\emph{U}(0) + \emph{n}) together with its post-lift shifted form.

\hypertarget{static-source-audit}{%
\subsubsection{Static source audit}\label{static-source-audit}}

A static audit scanned every .lean file and counted forbidden tokens.

\begin{longtable}[]{@{}ll@{}}
\toprule\noalign{}
Forbidden token & Count \\
\midrule\noalign{}
\endhead
\bottomrule\noalign{}
\endlastfoot
sorry & 0 \\
admit & 0 \\
axiom & 0 \\
opaque & 0 \\
\end{longtable}

This result is useful but is not a kernel result: lexical absence of
sorry does not establish that the files elaborate or that the proof
terms type-check. A lexical no-sorry / no-custom-axiom scan is not a
kernel proof.

\hypertarget{why-the-kernel-gate-is-not-closed}{%
\subsubsection{Why the kernel gate is not
closed}\label{why-the-kernel-gate-is-not-closed}}

The Lean toolchain could not be installed: the elan/Lean release-binary
download from github.com was blocked by the environment's egress policy
(HTTP 403; PyPI is allowlisted but GitHub release binaries are not), so
the following acceptance criteria remain unexecuted: a clean lake build;
kernel acceptance of every theorem body; a \#print axioms report for all
top-level declarations; a clean-environment independent reproduction;
and hash equality between two builds. The disposition is therefore
partially closed. The gap is narrowed to a mechanical execution: install
the pinned toolchain, run lake build, preserve stdout/stderr and append
the axiom report. Any elaboration failure must be repaired in the source
and rerun; it may not be converted into documentary proof closure.

\hypertarget{proposed-full-project-structure}{%
\subsubsection{Proposed full project
structure}\label{proposed-full-project-structure}}

FmiCns/ FSS/Basic.lean, FSS/Semantics.lean, FMI/Functions.lean,
FMI/SignalBoundary.lean, Projection/Conflation.lean,
Operations/SupportSchedule.lean, Operations/Nonreduction.lean,
Order/Composition.lean, Order/Span.lean, Order/Invariance.lean,
Fragmentation/Queue.lean, Fragmentation/Regional.lean, CNS/Lift.lean,
CNS/DynamicalBridge.lean, Governance/Options.lean,
Propagation/Regeneration.lean, GreatFilter/Constraints.lean,
CST/Paths.lean, CST/Transport.lean, CST/Closure.lean, Main.lean.

\hypertarget{declaration-skeletons}{%
\subsubsection{Declaration skeletons}\label{declaration-skeletons}}

structure FSS where Node : Type; Edge : Type; nodeType : Node
-\textgreater{} TypeTag; edgeType : Edge -\textgreater{} EdgeTag; metric
: Node -\textgreater{} Node -\textgreater{} Real; incident : Edge
-\textgreater{} Node -\textgreater{} Prop; boundary : Node
-\textgreater{} BoundaryRecord; admissible : (Node -\textgreater{} Node)
-\textgreater{} Prop.

structure Composition (X : FSS) where run : X.Node -\textgreater{}
X.Node; registered : Prop; reusable : Prop; spans : Prop. The order is
def CognitiveOrder (X : FSS) : Cardinal := Cardinal.mk \{c : Composition
X // c.registered /~c.reusable /~c.spans\}. The final implementation
should quotient by functional equivalence rather than count raw
implementations.

\hypertarget{top-level-theorem-signatures}{%
\subsubsection{Top-level theorem
signatures}\label{top-level-theorem-signatures}}

theorem firstOrderConflation \{W O Y : Type\} (q : W -\textgreater{} O)
(phi : W -\textgreater{} Y) \{x y : W\} (hq : q x = q y) (hphi : phi x
!= phi y) : Not (Exists fun decoder : O -\textgreater{} Y
=\textgreater{} decoder (q x) = phi x /~decoder (q y) = phi y).

theorem conditionalGreatFilter (core : CoreCnsPremises M) (physical :
PhysicalMapping M) (branch : BranchClosure M) (hazard :
NonSummableHazard M) (timescale : FilterTimescale M) (outcome :
TargetMatchingOutcome M) : PostIntelligenceFilter M. In the pinned
source package these top-level shapes carry proof terms for the bounded
core; the wider skeleton retains documentary placeholders that a
submission artifact must replace with kernel-checked proof terms and an
axiom report showing no sorryAx dependency.

\hypertarget{three-valued-reporting-outside-the-kernel}{%
\subsubsection{Three-valued reporting outside the
kernel}\label{three-valued-reporting-outside-the-kernel}}

Lean proves propositions or fails to construct a proof. Manuscript-level
reporting uses proved, disproved and undetermined. Undetermined is
assigned by the workflow when neither a proposition nor its negation is
certified within the declared formalization and budget; it is not a
third truth value inside Lean's logic.

\hypertarget{supplementary-note-12.-formalization-fidelity-semantic-equivalence-and-mutation-tests}{%
\subsection{Supplementary Note 12. Formalization-fidelity,
semantic-equivalence and mutation
tests}\label{supplementary-note-12.-formalization-fidelity-semantic-equivalence-and-mutation-tests}}

\emph{Local map.} This Note specifies the tests required before a Lean
proof can be treated as a faithful formalization of the manuscript
rather than a proof about a substituted object.

\hypertarget{formalization-fidelity-audit}{%
\subsubsection{Formalization-fidelity
audit}\label{formalization-fidelity-audit}}

Independent registrars receive frozen manuscript passages and construct
mappings to canonical declarations. Primary endpoints are exact recovery
of semantic identity, theorem dependencies, warrant status and
external-premise boundaries. Disagreement is resolved by typed
counterexample rather than majority vote.

\hypertarget{semantic-equivalence-tests}{%
\subsubsection{Semantic-equivalence
tests}\label{semantic-equivalence-tests}}

Required transformations include benign node renaming; coordinate
reparameterization preserving the metric; metric mutation preserving
labels; interaction mutation preserving distances; warrant mutation
preserving conclusions; and participant-fitness mutation preserving
aggregate mean. Expected verdicts are invariance under the first two and
semantic non-equivalence under target-relevant mutations.

\hypertarget{mutation-catalogue}{%
\subsubsection{Mutation catalogue}\label{mutation-catalogue}}

Replace cognitive order with Choice depth; reduce Register/Recall to
\emph{G}/\emph{L} without bridge types; swap support/schedule cells;
delete a theorem premise; promote a conditional claim to empirical fact;
collide aliases; remove a warrant edge; silently close an external
premise; remove a successor obligation; use an intervention label as
ground truth. Every mutation must either fail type checking, invalidate
a proof or produce an explicit undetermined dependency status.
Diagnostic tests built from registered invariants are controls; hidden
transfer is the primary evidence of propagation.

\hypertarget{supplementary-note-13.-executed-finite-computational-realization-and-gate-zero-audit}{%
\subsection{Supplementary Note 13. Executed finite computational
realization and Gate Zero
audit}\label{supplementary-note-13.-executed-finite-computational-realization-and-gate-zero-audit}}

\emph{Local map.} This Note reports the reproduced finite witness and
the positive v16 Gate Zero audit that regenerated it from the supplied
v8 lineage. It supplies constructive and conformance evidence but no
external AI or civilizational measurement.

The complete state was event sourced and updated by one transition
function. Each trajectory executed Register, Recall/activation, \emph{L}
or \emph{G} with optional lift/span, certificate validation, participant
Model/Evaluate/Decompose/Bridge/Governance, Store and reuse, Express and
consequence realization. The reported execution contained 12,800
trajectories and 102,400 events. These were reproduced bit-for-bit by
the positive v16 Gate Zero audit below.

\hypertarget{deterministic-witnesses}{%
\subsubsection{Deterministic witnesses}\label{deterministic-witnesses}}

The complete model was reported to pass 23 target-pair witnesses. Each
designated ablation failed or became correctly undetermined.

\begin{longtable}[]{@{}
  >{\raggedright\arraybackslash}p{(\columnwidth - 6\tabcolsep) * \real{0.2500}}
  >{\raggedright\arraybackslash}p{(\columnwidth - 6\tabcolsep) * \real{0.2500}}
  >{\raggedright\arraybackslash}p{(\columnwidth - 6\tabcolsep) * \real{0.2500}}
  >{\raggedright\arraybackslash}p{(\columnwidth - 6\tabcolsep) * \real{0.2500}}@{}}
\toprule\noalign{}
\begin{minipage}[b]{\linewidth}\raggedright
ID
\end{minipage} & \begin{minipage}[b]{\linewidth}\raggedright
Contract
\end{minipage} & \begin{minipage}[b]{\linewidth}\raggedright
Target distinction
\end{minipage} & \begin{minipage}[b]{\linewidth}\raggedright
Ablation result
\end{minipage} \\
\midrule\noalign{}
\endhead
\bottomrule\noalign{}
\endlastfoot
W01 & Register & external symbol becomes typed concept & fail \\
W02 & Express & stored verdict reconstructed externally & fail \\
W03 & Store & identity persists after context eviction & fail \\
W04 & Recall & stored bridge reactivates under later target & fail \\
W05 & L & local relation classification & fail \\
W06 & G & cross-region support preserves obligations & undetermined \\
W07-W09 & fitness coordinates & current, target and projected
distinctions & fail \\
W10-W15 & internal basis & action-conditioned consequence and correction
distinctions & fail / undetermined as specified \\
W16 & Promotion & promote admissible composite & undetermined \\
W17 & operand lift & separate lower-order-identical paired worlds &
undetermined \\
W18 & certified span & connect lifted operand to remote target & fail \\
W19 & Certificate & reject decoy bridge & fail \\
W20 & Governance & reject mean-positive participant harm & fail \\
W21 & carrier type & reject signal-space regress & fail \\
W22 & target revision & separate target movement from progress & fail \\
W23 & one generator & detect split transition semantics & fail \\
\end{longtable}

\hypertarget{higher-order-paired-worlds}{%
\subsubsection{Higher-order paired
worlds}\label{higher-order-paired-worlds}}

\begin{longtable}[]{@{}
  >{\raggedright\arraybackslash}p{(\columnwidth - 10\tabcolsep) * \real{0.1667}}
  >{\raggedright\arraybackslash}p{(\columnwidth - 10\tabcolsep) * \real{0.1667}}
  >{\raggedright\arraybackslash}p{(\columnwidth - 10\tabcolsep) * \real{0.1667}}
  >{\raggedright\arraybackslash}p{(\columnwidth - 10\tabcolsep) * \real{0.1667}}
  >{\raggedright\arraybackslash}p{(\columnwidth - 10\tabcolsep) * \real{0.1667}}
  >{\raggedright\arraybackslash}p{(\columnwidth - 10\tabcolsep) * \real{0.1667}}@{}}
\toprule\noalign{}
\begin{minipage}[b]{\linewidth}\raggedright
Condition
\end{minipage} & \begin{minipage}[b]{\linewidth}\raggedright
Target accuracy
\end{minipage} & \begin{minipage}[b]{\linewidth}\raggedright
Certified reach
\end{minipage} & \begin{minipage}[b]{\linewidth}\raggedright
Reuse
\end{minipage} & \begin{minipage}[b]{\linewidth}\raggedright
Containment
\end{minipage} & \begin{minipage}[b]{\linewidth}\raggedright
Harm
\end{minipage} \\
\midrule\noalign{}
\endhead
\bottomrule\noalign{}
\endlastfoot
local baseline & 0.311 & 0.000 & 0.000 & 1.000 & 0.000 \\
matched compute & 0.346 & 0.000 & 0.000 & 1.000 & 0.000 \\
matched memory & 0.316 & 0.000 & 0.000 & 1.000 & 0.000 \\
matched population & 0.333 & 0.000 & 0.000 & 1.000 & 0.000 \\
matched communication & 0.400 & 0.000 & 0.000 & 1.000 & 0.000 \\
certification only & 0.327 & 0.003 & 0.000 & 1.000 & 0.000 \\
lift only & 0.310 & 0.000 & 0.000 & 1.000 & 0.000 \\
span only & 0.317 & 0.000 & 0.000 & 1.000 & 0.000 \\
full CNS & 0.585 & 0.435 & 0.347 & 0.993 & 0.007 \\
residual-open CNS & 0.397 & 0.160 & 0.130 & 0.713 & 0.156 \\
\end{longtable}

\hypertarget{bounded-registration}{%
\subsubsection{Bounded registration}\label{bounded-registration}}

\begin{longtable}[]{@{}
  >{\raggedright\arraybackslash}p{(\columnwidth - 8\tabcolsep) * \real{0.2000}}
  >{\raggedright\arraybackslash}p{(\columnwidth - 8\tabcolsep) * \real{0.2000}}
  >{\raggedright\arraybackslash}p{(\columnwidth - 8\tabcolsep) * \real{0.2000}}
  >{\raggedright\arraybackslash}p{(\columnwidth - 8\tabcolsep) * \real{0.2000}}
  >{\raggedright\arraybackslash}p{(\columnwidth - 8\tabcolsep) * \real{0.2000}}@{}}
\toprule\noalign{}
\begin{minipage}[b]{\linewidth}\raggedright
Relation depth
\end{minipage} & \begin{minipage}[b]{\linewidth}\raggedright
Pipeline
\end{minipage} & \begin{minipage}[b]{\linewidth}\raggedright
Exact recovery
\end{minipage} & \begin{minipage}[b]{\linewidth}\raggedright
Relation retention
\end{minipage} & \begin{minipage}[b]{\linewidth}\raggedright
False confidence
\end{minipage} \\
\midrule\noalign{}
\endhead
\bottomrule\noalign{}
\endlastfoot
2 & generic summary & 0.300 & 0.561 & 0.333 \\
2 & relation extraction & 0.700 & 0.822 & 0.144 \\
2 & target-aware monitor & 0.867 & 0.928 & 0.000 \\
4 & target-aware monitor & 0.622 & 0.875 & 0.000 \\
6 & target-aware monitor & 0.233 & 0.750 & 0.000 \\
8 & target-aware monitor & 0.067 & 0.611 & 0.000 \\
12 & target-aware monitor & 0.000 & 0.249 & 0.000 \\
16 & target-aware monitor & 0.000 & 0.113 & 0.000 \\
\end{longtable}

A correlation-aware retention model achieved held-out log loss 0.258
versus 1.024 for homogeneous independence. The canonical codeword
reconstruction was identified at 0.940 under 1.5\% independent bit
flips.

\hypertarget{v16-gate-zero-audit-executed-closed-positive}{%
\subsubsection{v16 Gate Zero audit (executed,
closed-positive)}\label{v16-gate-zero-audit-executed-closed-positive}}

The Gate Zero question is whether the inherited version-8 numerical
results can be regenerated from the supplied frozen transition function,
source events, registry, configuration and conformance suite while
preserving the v16 definitions. The audit (identifier
EXP-GATE-ZERO-V16-EXECUTED-REPRODUCTION) obtained the v8 lineage (zip
SHA-256 f21b7ba6\ldots, matching the reference) and re-ran it.

\begin{longtable}[]{@{}ll@{}}
\toprule\noalign{}
Required artifact & Available \\
\midrule\noalign{}
\endhead
\bottomrule\noalign{}
\endlastfoot
v8 generator source & Yes (present, reproduced) \\
v8 event log & Yes (present, reproduced) \\
v8 trajectory table & Yes (present, reproduced) \\
v8 registry & Yes (present, reproduced) \\
v8 frozen configuration & Yes (present, reproduced) \\
v8 conformance suite & Yes (present, reproduced) \\
v8 environment lock & Yes (present, reproduced) \\
v8 source figures / source data & Yes (present, reproduced) \\
v8 checksums & Yes (present, reproduced) \\
v15 Article source & Yes \\
v15 Supplement source & Yes \\
v15 revision plan & Yes \\
\end{longtable}

The audit disposition is closed-positive. The execution manifest
validated (58 files, 0 hash mismatches); run\_experiment.py re-ran
cleanly (exit 0) in a clean venv pinned to numpy 2.3.5 / pandas 2.2.3 /
scipy 1.17.0 / networkx 3.6.1 / matplotlib 3.10.8 / scikit-learn 1.7.2
(plus tabulate, an unpinned optional report dependency); the conformance
tests passed; the model hash 797b16e3\ldots{} and registry hash
b7a6ef0e\ldots{} matched the manifest and reference; and all 18 data
files reproduced bit-for-bit after gzip decompression (including the 45
MB event log). Twelve key values regenerated at full precision: full-CNS
accuracy 0.5845588235294118, best matched comparator
0.39950980392156865, full-CNS certified reach 0.4349877450980392, best
comparator certified reach 0.003308823529411765, reuse
0.34681372549019607, Pareto containment 0.98671875, incoherent
containment 0.7171875, harmed fractions 0.012109375 / 0.155712890625,
lower-order decoder max accuracy 0.4805714285714286, canonical codeword
identification 0.94, and correlation-aware held-out log loss 0.258278.
The inherited numbers are therefore promoted from historical witness to
reproduced constructive witness (conformant under the single generator
and event lineage).

\hypertarget{supplementary-note-14.-primary-experiments-and-transfer-first-propagation}{%
\subsection{Supplementary Note 14. Primary experiments and
transfer-first
propagation}\label{supplementary-note-14.-primary-experiments-and-transfer-first-propagation}}

\emph{Local map.} This Note states the experiments that must be executed
after the formal objects are frozen. Transfer and regenerative
reconstruction are primary; conformance to built-in invariants is
diagnostic. Dispositions for this version are recorded per experiment.

\hypertarget{exp-gate-zero-v16-executed-closed-positive}{%
\subsubsection{EXP-GATE-ZERO-V16 (executed,
closed-positive)}\label{exp-gate-zero-v16-executed-closed-positive}}

Audit literal and executed definitions, label branches, surrogate
metrics, actual states, registry coverage, one-generator correspondence
and compatibility with the v16 semantic definitions. Executed for this
version with disposition closed-positive (Note 13): the supplied v8
lineage regenerated every inherited metric from source events under one
generator.

\hypertarget{exp-lean-formal-closure-partially-closed}{%
\subsubsection{EXP-LEAN-FORMAL-CLOSURE (partially
closed)}\label{exp-lean-formal-closure-partially-closed}}

Primary outcomes: clean build, zero sorry, no unregistered axioms,
complete theorem-manifest coverage, reproducible dependency report and
independent build. For this version a pinned source package exists and
passes a static no-sorry / no-custom-axiom scan, but no kernel build was
executed (Note 11); the disposition is partially closed.

\hypertarget{exp-formalization-fidelity-open}{%
\subsubsection{EXP-FORMALIZATION-FIDELITY
(open)}\label{exp-formalization-fidelity-open}}

Independent registrars map manuscript claims to canonical definitions
and theorem parameters. Primary endpoints are identity recovery,
dependency recovery, warrant classification and external-premise
separation.

\hypertarget{exp-semantic-equivalence-open}{%
\subsubsection{EXP-SEMANTIC-EQUIVALENCE
(open)}\label{exp-semantic-equivalence-open}}

Test benign and target-changing transformations from Note 12. Primary
endpoints are invariance sensitivity, mutation specificity and shortest
semantic witness.

\hypertarget{exp-lean-mutation-open}{%
\subsubsection{EXP-LEAN-MUTATION (open)}\label{exp-lean-mutation-open}}

Run the mutation catalogue. The primary endpoint is the fraction of
target-changing mutations that fail type checking or invalidate the
expected top-level theorem while benign transformations preserve it.

\hypertarget{exp-direct-recurrence-executed-closed}{%
\subsubsection{EXP-DIRECT-RECURRENCE (executed,
closed)}\label{exp-direct-recurrence-executed-closed}}

Generate a first growth interval, apply a certified order lift, then
continue growth under fixed order \emph{N}+1. The primary endpoint is a
preregistered change from increasing backlog to bounded residual
followed by renewed positive backlog drift. Executed for this version
across five regimes (Note 6, Table S-REC and Figures S-REC1-S-REC2); the
endpoint was met under fixed post-lift service and branch heterogeneity
and not met under proportional service, regenerative distribution or a
second order lift. The result is synthetic and model-relative.

\hypertarget{exp-lean-receiver-propagation-partially-closed}{%
\subsubsection{EXP-LEAN-RECEIVER-PROPAGATION (partially
closed)}\label{exp-lean-receiver-propagation-partially-closed}}

Conditions are prose-only, registry-only, direct typed-graph
formalization, CST/Lean package and corrupted controls. Receivers must
recover the dependency chain, distinguish proofs from external premises,
apply the formalization to hidden claim families, locate controlled
mutations and produce a package from which a second receiver succeeds
without author repair. Receiver and claim family are the units of
inference. For this version the protocol is frozen: six experimental
conditions, seven primary transfer tasks, eight primary endpoints
(including second-receiver regenerative success and author-repair
count), 14 hidden claim families instantiated as 28 randomized tasks, a
separately stored answer key, a scoring schema with fatal-error
categories, and a blinding and no-author-repair rule. The regenerative
endpoint is Regenerate(\emph{i}, \emph{j}; Γ, \emph{T}) = 1 only if
\emph{j} recovers a contract-preserving representation of Γ, executes it
on a hidden target family and produces a retransmission package from
which a second receiver succeeds without author repair; mere agreement
with the manuscript conclusion does not satisfy it. A bounded isolated
same-model receiver study was executed: two conditions (formal package
vs prose-only control), two Receiver A instances each in fresh contexts
seeing only their read-only packet, and one Receiver B regeneration from
a Receiver A transmission. All six core verdicts transferred in both
conditions (24/24); the registered theorem identity (T-CONFLATE) was
recovered by 2/2 formal-package receivers and 0/2 prose-only receivers;
the second-generation receiver reproduced all verdicts and the theorem
identity from the transmission alone with zero author repair. The
warranted claim is regenerative propagation across isolated same-model
instances; cross-model and human propagation remain open.

\hypertarget{receiver-hidden-family-inventory}{%
\subsubsection{Receiver hidden-family
inventory}\label{receiver-hidden-family-inventory}}

\begin{longtable}[]{@{}
  >{\raggedright\arraybackslash}p{(\columnwidth - 2\tabcolsep) * \real{0.5000}}
  >{\raggedright\arraybackslash}p{(\columnwidth - 2\tabcolsep) * \real{0.5000}}@{}}
\toprule\noalign{}
\begin{minipage}[b]{\linewidth}\raggedright
Family
\end{minipage} & \begin{minipage}[b]{\linewidth}\raggedright
Target distinction
\end{minipage} \\
\midrule\noalign{}
\endhead
\bottomrule\noalign{}
\endlastfoot
HF-CONFLATE & collapsed registration with target-different latent
states \\
HF-REFLECTION & reflective observability versus order lift \\
HF-REGISTER & boundary Register versus conceptual L/G endomorphism \\
HF-RECALL & memory Recall versus conceptual L/G endomorphism \\
HF-ORDER & novelty, span, reuse and non-equivalence \\
HF-RECURRENCE & fixed-service recurrence versus scaled-service
countermodel \\
HF-BRIDGE-LIFT-DYN & order lift plus contraction/closure/invariance \\
HF-BRIDGE-DYN-LIFT & new dynamical operator requiring a
prior-unavailable composition \\
HF-GOVERNANCE & participant completeness versus aggregate mean \\
HF-GF-TARGET & terminal versus observability target \\
HF-GF-PROMOTION & GF-3 conditional result versus GF-6 explanation \\
HF-SEM-RENAME & benign admissible isomorphism \\
HF-METRIC & label-preserving target-relevant metric mutation \\
HF-EXTERNAL & formal proof versus unmeasured external premise \\
\end{longtable}

\hypertarget{exp-cst-geometry-utility-open}{%
\subsubsection{EXP-CST-GEOMETRY-UTILITY
(open)}\label{exp-cst-geometry-utility-open}}

Compare direct typed graph formalization with the string/worldsheet
representation. CST geometry succeeds only if it improves encoding
invariance, semantic-drift detection, transfer, compression or receiver
reconstruction under matched resources. Failure demotes the
string-specific layer without weakening the core FSS/Lean program.

\hypertarget{supplementary-note-15.-registry-claim-ledger-and-manuscript-coverage}{%
\subsection{Supplementary Note 15. Registry, claim ledger and manuscript
coverage}\label{supplementary-note-15.-registry-claim-ledger-and-manuscript-coverage}}

\emph{Local map.} This Note specifies the source-of-truth records that
prevent prose, code and formal declarations from silently diverging.

Required registries are fss\_semantics\_v15, fmi\_functions\_v15,
operation\_topology\_v15, cognitive\_order\_v15,
fragmentation\_dynamics\_v15, cns\_bridge\_v15,
great\_filter\_constraints\_v15, cst\_semantic\_geometry\_v15,
lean\_theorem\_manifest\_v15, claim\_warrant\_ledger\_v15,
experiment\_registry\_v15 and version\_diff\_v14\_to\_v15. Every theorem
record contains premises, conclusion, proof certificate or proof
obligation, countermodels, dependencies, warrant, manuscript locations
and Lean declaration name. Every external-premise record states its
empirical interpretation, measurement protocol, status and the strongest
conclusion it can promote.

\hypertarget{principal-status-ledger}{%
\subsubsection{Principal status ledger}\label{principal-status-ledger}}

\begin{longtable}[]{@{}
  >{\raggedright\arraybackslash}p{(\columnwidth - 4\tabcolsep) * \real{0.3333}}
  >{\raggedright\arraybackslash}p{(\columnwidth - 4\tabcolsep) * \real{0.3333}}
  >{\raggedright\arraybackslash}p{(\columnwidth - 4\tabcolsep) * \real{0.3333}}@{}}
\toprule\noalign{}
\begin{minipage}[b]{\linewidth}\raggedright
Claim
\end{minipage} & \begin{minipage}[b]{\linewidth}\raggedright
Current v15 status
\end{minipage} & \begin{minipage}[b]{\linewidth}\raggedright
Promotion requirement
\end{minipage} \\
\midrule\noalign{}
\endhead
\bottomrule\noalign{}
\endlastfoot
Endogenous semantic identity & mathematical definition and prose theorem
& independent fidelity audit / executable formalization \\
T-CONFLATE & formal proof; source-realized in Lean core, kernel run
pending & kernel proof term \\
T-ORDER-INVARIANCE & prose proof under admissible-isomorphism premises &
Lean quotient implementation and build \\
Backlog divergence & formal proof & none beyond declared model \\
Regional fragmentation & conditional theorem & explicit obligation-graph
formalization \\
Direct recurrence & executed synthetic experiment (closed for the
declared model) & external instantiation for civilizational claims \\
T-CNS-LIFT & conditional theorem plus retained historical witness &
kernel closure and external recurrence evidence \\
Dynamical bridge & conditional theorem family & formal implementation of
published object \\
Option monotonicity & formal proof plus historical synthetic result &
measured error bounds \\
Regenerative propagation & executed bounded (same-model, second
generation, 0 author repairs) & cross-model / human receivers \\
GF necessary properties & decomposed conditional theorem & formal proof
terms / assumption challenge \\
FMI-CNS structural admissibility & GF-2 conditional theorem & physical
and branch mapping for GF-4/5 \\
Great Filter identification & open & hazard, timescale, residue and
comparative evidence \\
CST geometric utility & open & matched comparative experiment \\
Gate Zero conformance & executed, closed-positive (v8 lineage reproduced
bit-for-bit) & none beyond the declared witness class \\
Lean proof closure & source-complete, static-audit clean; kernel build
open & pinned clean build, axiom audit, independent reproduction \\
\end{longtable}

\hypertarget{supplementary-note-16.-external-mappings-standalone-completeness-and-completion-protocols}{%
\subsection{Supplementary Note 16. External mappings, standalone
completeness and completion
protocols}\label{supplementary-note-16.-external-mappings-standalone-completeness-and-completion-protocols}}

\emph{Local map.} This Note is the submission gate. It registers the
external-mapping certificate matrix, states what is already contained in
the package and states what must still be supplied before external
claims or Lean verification are promoted.

\hypertarget{external-mapping-certificate-matrix-registration-closed}{%
\subsubsection{External-mapping certificate matrix
(registration-closed)}\label{external-mapping-certificate-matrix-registration-closed}}

Each external mapping has a stable identity, formal object, minimum
evidence, strongest permitted claim and explicit defeater. Registration
closure prevents silent promotion; every mapping below remains open at
the instantiation level.

\begin{longtable}[]{@{}
  >{\raggedright\arraybackslash}p{(\columnwidth - 6\tabcolsep) * \real{0.2500}}
  >{\raggedright\arraybackslash}p{(\columnwidth - 6\tabcolsep) * \real{0.2500}}
  >{\raggedright\arraybackslash}p{(\columnwidth - 6\tabcolsep) * \real{0.2500}}
  >{\raggedright\arraybackslash}p{(\columnwidth - 6\tabcolsep) * \real{0.2500}}@{}}
\toprule\noalign{}
\begin{minipage}[b]{\linewidth}\raggedright
Mapping
\end{minipage} & \begin{minipage}[b]{\linewidth}\raggedright
Minimum evidence
\end{minipage} & \begin{minipage}[b]{\linewidth}\raggedright
Strongest permitted statement
\end{minipage} & \begin{minipage}[b]{\linewidth}\raggedright
Status
\end{minipage} \\
\midrule\noalign{}
\endhead
\bottomrule\noalign{}
\endlastfoot
Cross-FSS semantic principle & Target-specific invariant-preserving map
& The same structural principle can be instantiated across FSSs &
Open \\
Strong interaction ↔ Register & Support factorization and
schedule-commutation witnesses & Candidate physical realization &
Open \\
Weak interaction ↔ Recall & Support nonfactorization and precedence
witnesses & Candidate physical realization & Open \\
Electromagnetism ↔ L & Support factorization and precedence witnesses &
Candidate physical realization & Open \\
Gravity ↔ G & Support nonfactorization and commutation witnesses &
Candidate physical realization & Open \\
Current AI trace mapping & External version-locked collapsed-pair traces
with hidden targets & Testable external mapping & Open \\
Historical order lifts & Preregistered uniform classification and
independent evidence & Candidate historical instances & Open \\
Great Filter branch closure & Heterogeneous branch evidence or
defensible external model & Conditional branch premise & Open \\
Hazard non-summability & Longitudinal hazard and correction/dispersal
evidence & Conditional hazard premise & Open \\
Target residue suppression & Outcome- and timescale-specific residue
evidence & Conditional outcome premise & Open \\
\end{longtable}

The physical operation assignments are not standard measurements read
from established interaction ranges; their support and schedule
definitions are target-relative functional claims requiring explicit
factorization, nonfactorization, commutation and precedence
certificates. A mapping can now fail at a named certificate without
affecting the conceptual theorem package, while a successful certificate
promotes only the claim level it supports.

\hypertarget{standalone-completeness-checklist}{%
\subsubsection{Standalone-completeness
checklist}\label{standalone-completeness-checklist}}

\begin{longtable}[]{@{}
  >{\raggedright\arraybackslash}p{(\columnwidth - 2\tabcolsep) * \real{0.5000}}
  >{\raggedright\arraybackslash}p{(\columnwidth - 2\tabcolsep) * \real{0.5000}}@{}}
\toprule\noalign{}
\begin{minipage}[b]{\linewidth}\raggedright
Requirement
\end{minipage} & \begin{minipage}[b]{\linewidth}\raggedright
Disposition
\end{minipage} \\
\midrule\noalign{}
\endhead
\bottomrule\noalign{}
\endlastfoot
Typed metric FSS and endogenous semantics & installed \\
Reflective FMI and four signal residuals & installed \\
Canonical support and schedule definitions & installed \\
Register/Recall type nonreduction & prose proof installed \\
Intrinsic order and finite certificates separated & installed \\
Representation invariance & prose proof installed \\
Backlog and regional fragmentation separated & installed \\
CNS count lemma separated from substantive lift theorem & installed \\
Bridge to published dynamical CNS & installed conditionally \\
Great Filter lemmas decomposed & installed \\
CST metric separated from optional action & installed \\
Lean module and theorem manifest & installed; pinned source + static
audit installed; kernel build pending \\
v16 Gate Zero & executed; closed-positive (v8 lineage reproduced) \\
Direct recurrence & executed; closed for the declared synthetic model \\
External-mapping certificate matrix & installed (registration-closed) \\
Same-model receiver propagation & executed (bounded, isolated
same-model); cross-model/human pending \\
Repository and reviewer access & pending \\
\end{longtable}

\hypertarget{external-completion-protocols}{%
\subsubsection{External completion
protocols}\label{external-completion-protocols}}

\textbf{Lean build and independent validation.} Pin the toolchain,
replace every documentary sorry, print axiom dependencies, reproduce in
an independent environment and archive build logs.
\textbf{Formalization-fidelity audit.} Freeze the Article, Supplement,
registry and declarations; recruit independent registrars; score
identity, dependency and warrant recovery. \textbf{Receiver propagation
study.} Use fresh human or model receivers, hidden target families and
second-receiver retransmission without author repair; analyse receiver
and claim family, not individual prompts, as the units of inference.
\textbf{External trace study.} Use version-locked systems, matched
tasks, genuinely external traces and a target-aware three-valued
monitor; synthetic traces are not samples from current systems.
\textbf{Historical classification.} Freeze the lift rubric and
anti-tuning rules before selecting candidate innovations.
\textbf{Physical operation mapping.} Test minimal support, factorization
and commutation/precedence witnesses for the four physical assignments;
failure narrows the cross-FSS hypothesis without defeating the
conceptual theorems. \textbf{Great Filter branch and hazard studies.}
Test branch inheritance, regenerative thresholds, hazard
non-summability, correction/dispersal timescales and residue
suppression. \textbf{Repository closure.} Deposit code, data, registry,
Lean files, environment lockfiles, checksums, source figures and
manuscript source with reviewer access and a persistent DOI before
publication, including regeneration of the absent v8 executable lineage.

The minimum submit-ready projection is the core formal article, a
conformant finite witness (either the regenerated v8 lineage or the
executed recurrence result standing in its declared model), a clean Lean
build or a clearly narrower non-Lean formal positioning, and at least
one independent external anchor. The present v15 draft does not invent
those missing results.

\hypertarget{references}{%
\subsection{References}\label{references}}

\begin{enumerate}
\def\labelenumi{\arabic{enumi}.}
\tightlist
\item
  Gärdenfors, P. Conceptual Spaces: The Geometry of Thought (MIT Press,
  2000).\\
\item
  Collins, A. M. \& Loftus, E. F. A spreading-activation theory of
  semantic processing. Psychol. Rev.~82, 407-428 (1975).\\
\item
  Krenn, M. \& Zeilinger, A. Predicting research trends with semantic
  and neural networks with an application in quantum physics. Proc. Natl
  Acad. Sci. USA 117, 1910-1916 (2020).\\
\item
  Woolley, A. W., Chabris, C. F., Pentland, A., Hashmi, N. \& Malone, T.
  W. Evidence for a collective intelligence factor in the performance of
  human groups. Science 330, 686-688 (2010).\\
\item
  Hong, L. \& Page, S. E. Groups of diverse problem solvers can
  outperform groups of high-ability problem solvers. Proc. Natl Acad.
  Sci. USA 101, 16385-16389 (2004).\\
\item
  Lazer, D. \& Friedman, A. The network structure of exploration and
  exploitation. Admin. Sci. Q. 52, 667-694 (2007).\\
\item
  Lorenz, J., Rauhut, H., Schweitzer, F. \& Helbing, D. How social
  influence can undermine the wisdom of crowd effect. Proc. Natl Acad.
  Sci. USA 108, 9020-9025 (2011).\\
\item
  Murdoch, W. J., Singh, C., Kumbier, K., Abbasi-Asl, R. \& Yu, B.
  Definitions, methods, and applications in interpretable machine
  learning. Proc. Natl Acad. Sci. USA 116, 22071-22080 (2019).\\
\item
  Rudin, C. Stop explaining black box machine learning models for high
  stakes decisions and use interpretable models instead. Nat. Mach.
  Intell. 1, 206-215 (2019).\\
\item
  Christiano, P., Shlegeris, B. \& Amodei, D. Supervising strong
  learners by amplifying weak experts. Preprint at arXiv:1810.08575
  (2018).\\
\item
  Irving, G., Christiano, P. \& Amodei, D. AI safety via debate.
  Preprint at arXiv:1805.00899 (2018).\\
\item
  Burns, C. et al.~Weak-to-strong generalization: eliciting strong
  capabilities with weak supervision. Preprint at arXiv:2312.09390
  (2023).\\
\item
  Korbak, T. et al.~Chain of thought monitorability: a new and fragile
  opportunity for AI safety. Preprint at arXiv:2507.11473 (2025).\\
\item
  Emmons, S. et al.~When chain of thought is necessary, language models
  struggle to evade monitors. Preprint at arXiv:2507.05246 (2025).\\
\item
  Williams, A. E. Cognitive near-singularity: a possibility theorem and
  a witness-based protocol with a corpus case study. Front. Artif.
  Intell. 9, 1746633 (2026).\\
\item
  Wilson, E. B. Probable inference, the law of succession, and
  statistical inference. J. Am. Stat. Assoc. 22, 209-212 (1927).\\
\item
  Sandve, G. K., Nekrutenko, A., Taylor, J. \& Hovig, E. Ten simple
  rules for reproducible computational research. PLoS Comput. Biol. 9,
  e1003285 (2013).\\
\item
  Ćirković, M. M. The Great Silence: Science and Philosophy of Fermi's
  Paradox (Oxford Univ. Press, 2018).\\
\item
  Frank, A. \& Sullivan, W. T. A new empirical constraint on the
  prevalence of technological species in the universe. Astrobiology 16,
  359-362 (2016).\\
\item
  Wright, J. T. et al.~The G infrared search for extraterrestrial
  civilizations with large energy supplies. I. Background and
  justification. Astrophys. J. 792, 26 (2014).\\
\item
  Polyakov, A. M. Quantum geometry of bosonic strings. Phys. Lett. B
  103, 207-210 (1981).\\
\item
  Shannon, C. E. A mathematical theory of communication. Bell Syst.
  Tech. J. 27, 379-423, 623-656 (1948).\\
\item
  de Moura, L., Kong, S., Avigad, J., van Doorn, F. \& von Raumer, J.
  The Lean theorem prover (system description). In Automated Deduction -
  CADE-25, 378-388 (Springer, 2015).\\
\item
  Lean FRO. The Lean Language Reference, version 4.32.0 (2026).\\
\item
  Carneiro, M. Lean4Lean: towards a verified typechecker for Lean, in
  Lean. Preprint at arXiv:2403.14064 (2024).
\end{enumerate}
